\documentclass[11pt,letterpaper]{article}
\usepackage[T1]{fontenc}
\usepackage{lmodern,amsmath,amssymb,amsthm,booktabs,array,longtable}
\usepackage[margin=1in]{geometry}
\usepackage{microtype,xurl,enumitem}
\usepackage[dvipsnames]{xcolor}
\usepackage[colorlinks=true,linkcolor=MidnightBlue,urlcolor=MidnightBlue,citecolor=MidnightBlue]{hyperref}
\usepackage{fancyhdr}
\pagestyle{fancy}\fancyhf{}
\fancyhead[L]{\small PA consistency machines}
\fancyhead[R]{\small 13 November 2026}
\fancyfoot[C]{\thepage}
\setlength{\headheight}{14pt}
\setlength{\emergencystretch}{2em}
\setlist{topsep=4pt,itemsep=3pt}
\newtheorem{theorem}{Theorem}[section]
\newtheorem{proposition}[theorem]{Proposition}
\newcommand{\PA}{\mathrm{PA}}
\newcommand{\Con}{\operatorname{Con}}
\newcommand{\Bit}{\operatorname{Bit}}
\newcommand{\Prf}{\operatorname{Prf}}
\newcommand{\file}[1]{\texttt{\nolinkurl{#1}}}
\title{PA consistency machines:\\comparison, source compression, and proof obligations}
\author{Comparative construction and verification}
\date{13 November 2026}
\begin{document}
\maketitle
\begin{abstract}
We compare the incoming U208 constructions, James Chang's PA209 snapshot,
and the published PA196 machine. The nine designated U208 outputs form two
rooted transition-graph classes; their common contextual fold is not a set
of independent savings. We give a new source representation with 158
distinct internal nodes, compared with PA196's 159, preserving all 3,812
instructions and every ordinal control-flow edge. One representation has maximum depth
37 and maximum zero run six. A second, physically compiled representation
uses 196 working states: its six new contextual folds save a reader state,
but an additional return-scanner state is needed. This variant passes
all 4,959 arbitrary-value execution contracts and exact blank startup. We also give a direct
syntactic route to PA196's missing internally stated consistency direction,
and counterinterpretations showing that neither membership-rewriting head
is derivable from the remaining exact native rules. \textbf{No binary
Turing machine with fewer than 196 working states has been validated in
this investigation.} Source compression does not by itself achieve that
target. The packet supplies complete data and finite verifiers so that the
established results can be reused without relying on an optimizer's output
or confusing a candidate allocation with a correct machine.
\end{abstract}

\section{Scope and result status}

The requested endpoint is an explicit deterministic, single-tape,
two-symbol machine $M$, starting on an all-zero tape, with fewer than 196
working states and
\[
       \PA\vdash \neg\operatorname{Halt}(M)\leftrightarrow\Con(\PA).
\]
Working states exclude a distinguished halt control. We retain the binary
alphabet and blank-input convention of the compared constructions. State
labels need not be consecutive, so the largest label is not a state count.

The packet contains both new source representations and a validated new
196-state physical layout. PA196's reader states already implement several
pairs of source nodes through contextual sharing. Reparenthesizing the
source can destroy those savings or invalidate a tape invariant. A nominal
195-state allocation fails a concrete register-frame test; repairing it
requires one additional scanner state. The rejected table is clearly
labelled and is not presented as a new record.

There are three different levels of evidence in this packet:
\begin{enumerate}
\item Exact finite checks: source/control-flow equality, machine census,
transition-graph comparison, startup replay, and complete-tape rejection
of tempting merges; also symbolic induction certificates covering all
register values and finite spare lists for the new 196-state variant.
\item Mathematical arguments: soundness-preserving source identity,
the two independence counterinterpretations, and the direct syntactic
consistency proof construction in the accompanying note.
\item Outstanding implementation work: a sub-196 physical allocation,
and a newly exported
kernel-checked PA theorem for the reverse consistency implication.
\end{enumerate}
The finite experiments neither prove nonhalting by simulation nor establish
a lower bound on all possible 195-state machines.

\section{Pinned sources and comparison}

The PA196 source is pinned to
\file{jamesyc/BusyBeaver@407cdc356f6b586635e71ac92be394255f4bfea4};
the requested PA209 source is pinned to
\file{69d98241a721ead20b4f0b17df9062950b40fbfe}.
The incoming archives and the additional comparison/search work are pinned
to \file{VladimirReshetnikov/ProveIt@3a7ddf022dc41b77db860f8db3f888002ddc4f47}
\cite{pa196,pa209,incoming,comparison}.

\begin{table}[ht]
\centering
\begin{tabular}{@{}lrrrr@{}}
\toprule
Construction & Working states & Rows & Source sites & Startup steps\\
\midrule
PA209 &209&418&3,948&58,210\\
U208 families &208&416&3,909&48,140\\
PA196 &196&392&3,812&58,154\\
New validated variant &196&392&3,812&59,164\\
\bottomrule
\end{tabular}
\caption{The U208 source count is the reported count for the reviewed
families. Startup counts for the predecessor and U208 are from the pinned
sources; both 196-state startups are independently replayed by this packet.}
\end{table}

\subsection{What is actually proved in the upstream formal developments}

Both James Chang snapshots explicitly export the internal implication
\[
        \PA\vdash\neg H\longrightarrow\Con(\PA),
\]
where $H$ is the exact arithmetic halting sentence for the literal table.
Their external nonhalting proofs and standard-model soundness results
support the advertised independence conclusions. Those declarations are
not themselves a formal export of the biconditional requested here
\cite{pa196result,pa209result}.

This is a statement about the theorem inventory, not an assertion that the
reverse implication is false. The note
\file{notes/internal_soundness.tex} develops a direct syntactic construction
for it. In particular, external validity in a hereditarily finite model is
not silently substituted for a uniform PA proof about nonstandard proof
codes and nonstandard execution lengths.

\subsection{The U208 archives contain two machine families}

All eleven incoming ZIP identities were checked against their Git blob
identities. Five construction packages designate nine final outputs.
Canonically numbering the rooted, action-labelled transition graphs gives
two classes, with five and four outputs respectively. There are seven
different native byte strings and six native-labelled transition functions.
These distinctions account for comments, row order, and state names; they
are not nine independent mathematical constructions.

\begin{table}[ht]
\centering\small
\begin{tabular}{@{}ll@{}}
\toprule
Class A & Class B\\
\midrule
ConPA primary & Conditional compatibility\\
Conditional primary & Entry compatibility\\
Entry primary & History compatibility\\
History primary & Research primary\\
Research plain & \\
\bottomrule
\end{tabular}
\caption{The nine designated final U208 outputs. The exact archived paths,
hashes, and copied table names are in \file{data/u208_inventory.json}.}
\end{table}

The remaining incoming packages are research dossiers, comparison work,
source specializations, and a unary-theory calibration. They were included
in the investigation, but were not counted as extra successful transition
tables. The historical comparison verifier was also replayed: its 3,744
transition-conjugacy equations, 48,140-step startup calculations, manifest
checks, and two negative controls passed. This verifies those finite
claims; it is not a new end-to-end PA formalization.

\subsection{Why the U208 fold does not simply subtract another state}

The U208 constructions perform the same contextual fold of native controls
172 and 211. Equivalent presentations retain different names for the merged
control. The proof uses information about recently written tape cells and
the incoming-state census. In the Joint/Plain case the merged zero branches
take seven and three child steps; the one branches take one. At local
boundaries the whole tape is unchanged. Thus the corresponding times satisfy
\[
 t'(n)=n+6N_{(172,0)}(n)+2N_{(211,0)}(n),\qquad n\le t'(n)\le7n.
\]
This is a useful PA-formalizable local equivalence. It applies to that
parent table and its verified context invariant. It cannot be subtracted
again merely because several archives describe it.

The exact mixed-entry pattern used by that fold is absent from PA196.
The newer ProveIt comparison already records this fact and substantial
unsuccessful local searches \cite{comparison}. Consequently a successful
combination requires a different contextual certificate or a changed source
layout; matching state numbers across the two families has no mathematical
meaning.

\subsection{What PA196 changes relative to PA209}

PA196's reduction is largely a change in the emitted calculus and its
representation, not the U208 fold in isolation. Its universal constructor is
\[
 U(a,B)=4p(B,a)+2,\qquad p(u,v)=\binom{u+v+1}{2}+v,
\]
so the pairing arguments are transposed. Implication, membership, and
equality retain residues $0,1,3$. The matching decoder changes as well.
The emitter removes the old seventh equality head, inserts guarded equality
vacuity into the separation gate, rearranges heads and antecedents, and
uses six contextual pairs in its AX body. Two additional obsolete routing
states are removed. The published census is 36 engine states, 13
leaf-routing states and 147 reader states \cite{pa196}.

\section{An exact 158-node representation of the PA196 source}

\subsection{The finite object}

The full candidate is \file{data/source158.json}; its SHA-256 is
\begin{quote}\small
\nolinkurl{7732de4da7b4c988046c77d9d7d2a6f6dbb7ea5c54a17901190add5b5e5043f0}.
\end{quote}
Each row contains a binary path, an instruction opcode, and an operand.
The paths are the leaves of a full binary tree in ordinal order. Equal
subtrees are stored once when counting the DAG. The five instruction types
are increment, conditional decrement, forward break, backward continue,
and halt. Nine logical registers are used.

\begin{table}[ht]
\centering
\begin{tabular}{@{}lrr@{}}
\toprule
Quantity & Published PA196 source & New representation\\
\midrule
Distinct internal nodes &159&158\\
Instruction occurrences &3,812&3,812\\
Maximum path length &38&37\\
Maximum consecutive zero run &6&6\\
Different resolved instruction rows &\multicolumn{2}{c}{0}\\
\bottomrule
\end{tabular}
\caption{Exact counts independently checked by \file{code/verify.py}.}
\end{table}

\begin{theorem}[Exact source identity]\label{thm:source}
Number the instruction occurrences of both source trees from $0$ to
$3811$. At every ordinal, the operation, its register operand if any,
and every resolved successor ordinal agree. Therefore, for every initial
register vector and the same ordinal control, the two source programs
have identical register/control executions, including halting.
\end{theorem}
\begin{proof}
The verifier parses the pinned original Lean tree, independently checks it
against \file{data/source196.json}, and reconstructs the new full binary
tree from its paths. It compares the complete ordered instruction words.
It then resolves both programs twice. One resolver uses subtree intervals
and the relevant ancestor edge; the other uses binary increment, suffix
cut, and zero extension to locate the successor leaf. All four 3,812-row
graphs agree. No unresolved root exits occur.

For an increment the identical next row and register operation give
identical successors. For a decrement, the same register is tested and
both the zero and positive successors agree. The two jump forms have
identical ordinal targets. Halt occurs at the same ordinal. Induction on
the number of source steps proves the claim. This is an identity of source
semantics for arbitrary registers, stronger than agreement on a sampled
prefix from the clean seed.
\end{proof}

The clean seed remains ordinal 25 and register vector
$(2,0,0,0,0,0,0,0,0)$. Theorem~\ref{thm:source} does not transfer the
literal tape encoding of that seed to a different physical machine.
The separately compiled variant below discharges those obligations for
its own source representation and literal table.

\subsection{How the representation was found}

The search treats the straight-line AX emitter as a word in existing
macro blocks and single accumulator increments. It considers repeated
substrings as candidate grammar phrases and solves a finite mixed-integer
segmentation problem. Activating a phrase and choosing a path through its
positions gives a grammar; a phrase split into $k$ children costs $k-1$
binary grouping nodes. Balanced expansion, followed by semantics-preserving
reassociation, supplies a binary source tree.

The one-node gain exploits reuse across the emitter/macro boundary:
the pair of adjacent accumulator increments already occurs as
\file{node60} in the multiplication-by-four component of the pairing macro.
An AX grammar that reuses that exact subtree can save a full-program DAG
node even when its local grammar objective alone does not decrease.
Subsequent reassociation restores the six-zero bound and improves maximum
depth. The final data and the exact source identity are independent of
solver settings or the randomized discovery path.

This search is broader than the earlier bounded one-, two-, and selected
three-rotation screens in ProveIt. Its result is an explicit improvement
over their 159-node source baseline. It is not a claim that 158 is optimal.

\subsection{Why a source-node saving does not automatically save a state}

The equation ``159 source nodes become 158, therefore 196 states become
195'' is invalid. In PA196, several distinct nodes share a reader state
because their incoming tape contexts distinguish them. A new tree may
retain fewer such compatible pairs or require extra helper states. The
ordinary 158-node source does not preserve all six old AX folds.

The physical allocator must satisfy all of the following jointly: dispatch
to the intended instruction, correct arbitrary-register tape behavior,
preservation of the register/header layout, and a correct blank-tape
startup. A successful source-prefix check is insufficient. Some decoder
states are revisited during primitive execution, so treating their rows
as freely writable can corrupt a register operation. The next section
gives a concrete example of exactly this failure mode.

\section{A validated 196-state regrouping and its rejected predecessor}

A joint source-and-fold search produces a different 158-internal-node
representation, supplied as \file{data/variant196_source.json}. It retains
the same 3,812 ordered primitive instructions and ordinal control-flow
edges, has depth 36, and has maximum zero run seven. Six compatible
contextual folds give a nominal 195-state table. The entry to the shared
copy-guard reader is constrained to retain its incoming one bit.

The nominal 195-state table is nevertheless wrong as a source interpreter.
At source ordinal 1003, a zero decrement of register 1 begins with seven
trailing zero address bits. With gap 48, nine zero registers and spare
list $[0]$, the machine first misroots after 81 steps. After 276 steps it
reaches a clean root with gap 49 and spare list $[0,1]$, rather than the
specified unchanged file and gap. The seven-zero return scanner has
mistaken a path region for the return delimiter. The complete tape
counterexample is \file{data/rejected195_example.json}; it is replayable by
\file{code/check_concrete.py}. This refutes the intended first-return-to-clean-root instruction contract,
not an eventual blank-tape halting equivalence by itself. With register 1
initially equal to 1, the same candidate reaches a wrong clean root after
278 steps: it leaves register 1 equal to 1 rather than decrementing it to
0, and takes the zero-branch successor. This also violates the logical
register effect, independently of the padding and spare-list discrepancy.

The repair changes native row $34/0$ from $(0,L,75)$ to $(0,L,2000)$
and adds
\[
 \delta(2000,0)=(0,L,75),\qquad
 \delta(2000,1)=(1,L,2).
\]
The new scanner distinguishes eight consecutive zeros. The repaired
machine has 196 working states and is supplied in both native and dense
forms, \file{data/variant196.native.tm} and \file{data/variant196.tm}.
In dense form its start is 0, its working states are $0,\ldots,195$,
and its separate halt is 196. The native table's SHA-256 is
\begin{quote}\small
\nolinkurl{f6594edc95efed02f26b88f71c9b1f7699965a258739836f24e6b641887b0cd9}.
\end{quote}

All 4,959 symbolic instruction cases pass with gap $48+g$, for arbitrary
$g\in\mathbb N$, nine independent natural register values, an arbitrary
first spare value, and an arbitrary finite remaining spare list. The
certificates use 1,177,411 symbolic inferences and cover all 392 table
rows. Their durations have nonnegative affine coefficients and positive
constant terms; the smallest constant is 3. Unary self-loop induction
and four rederived one-block traversal lemmas handle unbounded register
values and lists. These are executable mathematical certificates, not
finite small-register sampling and not proof-assistant kernel certificates.

From the blank tape, the repaired machine reaches exactly the baseline
clean frame after 59,164 steps: native root 84, head $-2$, the same 66
one-cells, gap 48, source ordinal 25, registers $(2,0,\ldots,0)$, and 26
spare blocks of value 1. Startup and runtime are checked separately.
The standalone note \file{notes/regrouped196_validation.tex} explains
certificate soundness and the PA-formalizable halting transfer in detail.

\begin{theorem}[Halting transfer for the repaired variant]
Let $R$ be the repaired table and $M_{196}$ the pinned baseline. The exact
source identity, startup checks and uniform positive runtime contracts,
with the elementary arithmetic soundness proof of their certificate rules,
give the paper-level conclusion
\[
 \PA\vdash \operatorname{Halt}(R)
             \leftrightarrow\operatorname{Halt}(M_{196}).
\]
Consequently the consistency conclusions for the baseline transfer to $R$.
\end{theorem}
\begin{proof}
A finite source halt gives a physical halt by concatenating startup and
primitive contracts. Conversely, if a literal halt occurred at time $T$
but the source did not halt, $T+1$ source steps and their positive physical
segments would give a literal trace longer than $T$ with a working
endpoint. Determinism contradicts the missing outgoing transition at the
alleged halt. These arguments are finite-sequence inductions expressible
in PA. Apply them to both implementations of the identical ordinal source.
\end{proof}

The new layout is a successful intermediate result at the existing bound.
Its extra scanner state consumes the reader-state saving; it does not
satisfy the requested strict inequality.

\section{A complete-tape rejection of a plausible merge}

One near miss in the existing local searches identifies dense states 96
and 103, corresponding to native controls 109 and 124. It is tempting
because local representative contexts look compatible. We tested the
actual full root-frame primitive contract.

Take source ordinal 16, path \texttt{0100110110}, whose instruction is
increment logical register 6. Use header gap 56, all nine registers zero,
and spare list $[0]$. Redirect native 109 to 124 and otherwise retain the
table. The exact results are:
\begin{center}
\begin{tabular}{@{}lrr@{}}
\toprule
Quantity & PA196 & Merged candidate\\
\midrule
Steps to the displayed successor frame &596&574\\
Successor path &\texttt{0100110111}&\texttt{0100110111}\\
Incremented register &6&8\\
Final spare list &$[0,1]$&$[0,1]$\\
\bottomrule
\end{tabular}
\end{center}
The complete tapes are checked, including every register block. Thus a
correct program location does not imply a correct implementation of the
source instruction. The script \file{code/reject_merge.py} reproduces this
finite counterexample. It disproves the intended first-return root-frame contract, not the
candidate's eventual blank-input halting equivalence.

Another tested shortcut shortened the seven-zero return scan to six by
deleting native state 34. Even though six-zero runs in the source paths
are trailing, transient tapes also matter. During startup the shortened
scan first returns to a different root position on the 54th root visit:
the original head is at $-2$, the altered head at $-1$, on identical tapes.
This is enough to reject the proposed identity simulation; no nonhalting
claim is drawn from a finite run.

\section{A syntactic route to the internal consistency direction}

The accompanying seven-section note
\file{notes/internal_soundness.tex} develops the following route to
\[
        \PA\vdash\Con(\PA)\to\neg H_{196}.
\]
It is a paper-level proof construction, not a newly kernel-checked Lean
export. Together with the published internal implication it supplies a
mathematical route to the full biconditional for the \emph{196-state}
baseline and, by the preceding transfer, the repaired 196-state variant.
Neither table is below the requested state bound.

Translate numeric formulas directly into PA, retaining source variable
labels:
\[
 M(a,b)\mapsto\Bit(x_a,x_b),\quad E(a,b)\mapsto x_a=x_b,
 \quad U(a,B)\mapsto\forall x_a\,\tau(B).
\]
The source quantifier's reversed pairing is decoded in the correct order.
Keeping labels avoids the binder/property-code mismatch that arises when
recoding through the older U208 formula convention. Code zero is truth;
the exceptional equality $I(0,0)=0$ is handled by a fixed logical
equivalence. Code one is self-membership, and PA uniformly proves
$\neg\Bit(z,z)$, so it serves as falsity under every assignment.

The nonlogical core is a uniform proof compiler for finite bit
comprehension. For an arbitrary input formula $R(z,\vec a)$, instantiate
full PA induction with
\[
 Q(n):=\exists s<2^n\ \forall z<n\,
                  \bigl(\Bit(s,z)\leftrightarrow R(z,\vec a)\bigr).
\]
At the successor step, retain $s$ when $R(n,\vec a)$ is false and choose
$s+2^n$ when it is true. This uses excluded middle, not a decision procedure
for $R$. The proof is generated by fixed syntactic templates, one induction
instance, freshening and substitution. It therefore does not define an
arithmetic truth predicate.

For the guarded separation gate, use bound $x_a+x_b+1$ and the selected
property
\[
 (\Bit(x_a,z)\vee z=x_b)\wedge\tau(c)[z/x_b].
\]
The case $a=b$ is treated separately because quantification then binds
both occurrences in the membership antecedent. A parallel stack retains
PA proofs of all live formula codes; a successful Shift halt test produces
a PA refutation.

The cleanest endpoint argument works inside PA under the explicit
hypothesis $\Con(\PA)$. That proof-witness invariant keeps every command
live. The existing all-model positive scheduler and physical simulation
contracts then build literal traces of arbitrarily large length with
working endpoints. Determinism excludes a finite halt. The note identifies
the scheduler contracts used and avoids confusing an absorbing coded
source halt with the literal target's missing outgoing transition.

\section{Why neither membership head can simply be deleted}

The two heads, up to exchanging antecedents, are
\begin{align*}
 L(a,b,d)&:\ E(a,b)\to(M(a,d)\to M(b,d)),\\
 R(a,b,d)&:\ E(a,b)\to(M(d,a)\to M(d,b)).
\end{align*}
Deleting one and deriving it from pointed separation would shorten the
emitter. The note \file{notes/membership_independence.tex} proves that this
particular justification cannot work for the unchanged native rules.

Let $\sigma$ swap 2 and 3 and fix all other naturals. Put $f_1=g_1=\sigma$
and $f_i=g_i=\mathrm{id}$ otherwise. Interpret equality as actual equality.
To refute $L$, interpret
\[
           M(a,b):=\Bit(f_a(x_a),x_b).
\]
Every remaining head, rule and the complete guarded gate is valid, but
$M(1,2)\to(E(1,0)\to M(0,2))$ is false when
$x_1=x_0=2$ and $x_2=0$. To refute $R$, instead interpret
\[
           M(a,b):=\Bit(x_a,g_b(x_b)).
\]
Then $E(0,1)\to(M(2,0)\to M(2,1))$ is false when
$x_0=x_1=2$ and $x_2=4$, while all remaining rules and the full gate
are valid. The label dependence is legitimate: the raw native calculus
does not postulate arbitrary variable-substitution rules for its atoms.

The gate verification is the substantive part. In the first interpretation,
form the finite selected bitset and apply $f_c^{-1}$ to obtain the witness
at its fresh binder. In the second, use the images of the selected indices
under the bijection $g_b$. The identities
$\neg\Bit(\sigma(t),t)$ and $\neg\Bit(t,\sigma(t))$ handle the diagonal
case $a=b$. These constructions work for arbitrary decoded properties,
not only a bounded collection of examples.

The finite checker additionally verifies the displayed counterexamples,
27,648 surviving-head instances and 225,000 selected-set constructions.
Those samples supplement the universal mathematical proof. The result
does \emph{not} exclude a different PA interpretation of a weaker
one-head calculus and does not imply a Turing-machine state lower bound.

\section{Relation to ConciseIPRA and the remaining optimization problem}

The ConciseIPRA material concerns presentations of intuitionistic primitive
recursive arithmetic \cite{ipra}. Its useful design lesson here is to
make the syntactic kernel explicit and transport proofs through concrete
local rules. Its consistency strength cannot simply replace full PA in
this problem. Similarly, the byte size of a Python proof checker does not
measure the number of states in its binary, blank-tape implementation.

For the present source, the useful objective is the size of a \emph{complete
physical allocation}: source nodes, retained contextual identifications,
new helpers, routing, and engine interactions. Minimizing just source
length or just DAG nodes can make that objective worse. The 158-node
representations and the fully checked 196-state variant are explicit
starting points for that combined problem. They have not yet crossed
the physical 196-state bound.

\subsection{Head-order and antecedent variants}

A separate exact symbolic check covers all $6!=720$ orders of the six
nonpropositional AX heads. Independently exchanging the first two
antecedents of heads 3, 4 and 6 gives $720\cdot 8=5,760$ variants.
The fixed generalization and propositional heads, saved top, guarded gate,
parameter reads and cache clearing remain unchanged. The note
\file{notes/axiom_head_permutation.tex} proves that the existing total
stack macros remain valid and that replacing each head-selector position
by its inverse-permutation position gives a primitive recursive all-model
refutation compiler. Antecedent exchange uses the existing propositional
head and modus ponens. Thus the logical change is justified before any
physical savings are claimed.

A greedy grammar screen over all 5,760 words, followed by eight bounded
mixed-integer searches on selected words, did not yield a validated
sub-196 table. The best of those eight source grammars has 158 internal
nodes but does not immediately supply a geometry-safe physical allocation
with the required contextual folds. The solver gaps remain nonzero. These are constructive
upper bounds and restricted unsuccessful searches, not lower bounds or
optimality claims.

\section{Reproduction and packet contents}

Run from the extracted packet root:
\begin{quote}\ttfamily
python3 code/verify.py\\
python3 code/reject\_merge.py\\
python3 code/reject195.py\\
python3 code/check\_variant.py\\
python3 code/check\_membership\_countermodels.py
\end{quote}
The principal verifier needs Python's standard library only. It reparses
the original source, checks the two source graphs independently, verifies
the table counts and U208 classes, and replays all 58,154 startup steps
against the exact complete tape. Its output is reproducible and is recorded
in \file{data/verification.json}.
For the new 196-state table, separately replay:
\begin{quote}\small\ttfamily
python3 code/check\_startup.py --machine data/variant196.native.tm\par
\hspace*{1em}--output certificates/variant196/startup.json\par\medskip
python3 code/check\_runtime.py --machine data/variant196.native.tm\par
\hspace*{1em}--source data/variant196\_source.json --gap 48 --states 196\par
\hspace*{1em}--output certificates/variant196
\end{quote}
Each displayed wrapped command is one shell invocation. Do not disable
Python assertions with \file{-O} or \file{PYTHONOPTIMIZE}. The checkers
return a nonzero exit code on failure. The runtime replay normally takes
about one to two minutes in the environment used for this investigation.

The packet contains the original PA196 and PA209 tables, all nine designated
U208 native tables, exact source data, provenance and licenses, the
independent verifiers, and five LaTeX/PDF documents. The supporting notes contain the longer logical and runtime proofs.
The included 195-state table is explicitly rejected by a concrete
contract failure; only the repaired 196-state variant is validated.
The alphabet, blank-input convention, state-count convention and requested
logical strength have not been weakened.

\begin{thebibliography}{9}
\bibitem{pa196} J. Chang. \emph{Nonhalting and PA-independence for a
196-state Turing machine}. Pinned revision
\href{https://github.com/jamesyc/BusyBeaver/blob/407cdc356f6b586635e71ac92be394255f4bfea4/paper/pa196.tex}{407cdc356f6b586635e71ac92be394255f4bfea4}.
\bibitem{pa196result} J. Chang. \emph{PA196 formal theorem statements}.
\href{https://github.com/jamesyc/BusyBeaver/blob/407cdc356f6b586635e71ac92be394255f4bfea4/docs/result.md}{Pinned theorem inventory}.
\bibitem{pa209} J. Chang. \emph{Nonhalting and PA-independence for a
209-state Turing machine}.
\href{https://github.com/jamesyc/BusyBeaver/blob/69d98241a721ead20b4f0b17df9062950b40fbfe/paper/pa209.tex}{Requested PA209 snapshot}.
\bibitem{pa209result} J. Chang. \emph{PA209 formal theorem statements}.
\href{https://github.com/jamesyc/BusyBeaver/blob/69d98241a721ead20b4f0b17df9062950b40fbfe/docs/result.md}{Pinned theorem inventory}.
\bibitem{incoming} V. Reshetnikov, ProveIt repository. \emph{Incoming
consistency-machine archives}.
\href{https://github.com/VladimirReshetnikov/ProveIt/tree/3a7ddf022dc41b77db860f8db3f888002ddc4f47/docs/incoming}{Pinned archive collection}.
\bibitem{comparison} ProveIt. \emph{PA196 and the incoming U208 construction},
with associated source, allocation and local-fold experiments.
\href{https://github.com/VladimirReshetnikov/ProveIt/tree/3a7ddf022dc41b77db860f8db3f888002ddc4f47/Computability/BusyBeaver/Research/PA/pa196-vs-u208}{Pinned comparison subtree}.
\bibitem{ipra} ProveIt. \emph{ConciseIPRA}.
\href{https://github.com/VladimirReshetnikov/ProveIt/tree/3a7ddf022dc41b77db860f8db3f888002ddc4f47/Logic/PrimitiveRecursiveArithmetic/ConciseIPRA}{Pinned proposals and formal developments}.
\end{thebibliography}
\end{document}
